% Copyright 2026 by Robert Hildebrand
% CC BY-SA 4.0
\chapter{Project: Optimization in Machine Learning}
\label{ch:project-ml}
\todoChapter{Skeleton (2026-07-31). Choose dataset(s), write starter notebooks, decide depth of the LLM/neural-network training discussion.}

\section*{The problem}
Training a machine learning model \emph{is} an optimization problem —
usually a huge, nonconvex one attacked with the unconstrained methods of
Chapters~17--18 at massive scale. This project has students train models
themselves, treat the training process as an optimization experiment, and
then close the loop: use a trained model's predictions inside a
downstream optimization model (predict-then-optimize).

\section*{Learning goals}
Empirical risk minimization as nonlinear optimization; gradient descent,
stochastic gradients, momentum, and Adam as algorithmic refinements of
Chapter~18; why nonconvexity is survivable in practice; regularization as
a modeling choice; the scale of modern training (what it means that an LLM
is trained by SGD on billions of parameters); hyperparameter search as a
black-box optimization / heuristic problem; combining a learned demand
forecast with an inventory or routing model.

\section*{Stages}
\begin{enumerate}
  \item \textbf{By hand.} Logistic regression on a real dataset with your
    own gradient descent; verify against a library. Plot loss curves for
    step-size choices — connect to line search and convergence theory.
  \item \textbf{Neural network.} A small feed-forward network (or
    fine-tune a small pretrained model) with SGD/Adam; study batch size,
    learning-rate schedules, and stochasticity as a heuristic device.
  \item \textbf{Hyperparameters.} Random search vs.\ a simple simulated
    annealing over the configuration space — heuristics from Chapter~16
    applied to learning.
  \item \textbf{Predict, then optimize.} Feed the model's predictions
    (e.g.\ demand) into a MIP from Part~III; quantify how prediction error
    propagates into decision cost.
  \item \textbf{Report.} An experiment log with the discipline of an
    optimization study: baselines, ablations, and honest error bars.
\end{enumerate}

\section*{Deliverables and grading}
% TODO: rubric; dataset shortlist; compute guidance (all stages must run
% on a laptop or free-tier GPU).
